\chapter{大规模AI基础设施设计}
\label{ch:ai-infrastructure}

\section{引言}

随着生成式人工智能（Generative AI）的迅猛发展，构建支撑大规模模型训练的基础设施已成为数据通信领域面临的核心挑战之一。Meta在其2024至2026年的基础设施演进过程中，逐步构建了从两万四千块GPU到吉瓦（Gigawatt）级别AI集群的完整技术体系。这一演进过程不仅体现了计算规模的指数级增长，更揭示了传统网络架构在面对AI训练工作负载时的根本性局限。

AI训练工作负载与传统云计算流量存在本质差异。训练任务通常涉及数千乃至数万个加速器之间的密集通信，采用\textbf{RDMA（Remote Direct Memory Access，远程直接内存访问）}技术——这是一种允许网卡绕过操作系统内核、直接读写远程服务器内存的技术——通过UDP协议交换数据。这种通信模式产生了所谓的"\textbf{大象流}（elephant flows）"——持续时间极长、数据量巨大的流量。可以将其类比为高速公路上行驶的超长货车队：想象一下，一条高速公路上既有普通私家车，也有绵延数公里的超长重型卡车编队。传统路由策略像随机分配车道，如果多个超长车队碰巧被分配到同一条车道，这条车道就会彻底堵死，而其他车道却空空如也。

与此同时，AI训练还呈现出\textbf{低熵（low entropy）流量}特征——即由于参与集合操作的GPU数量有限，导致IP流的数量较少、模式重复可预测。这就好比在少数几条车道上承载大量车辆，即使总容量充足，局部拥堵也难以避免。这种低熵特性会引发哈希碰撞和链路拥塞。传统的基于哈希的\textbf{ECMP（Equal-Cost Multi-Path，等价多路径）}路由策略——这是一种基于哈希算法将流量分散到多条等价路径的路由策略——在这种场景下表现不佳，因为流量无法均匀分布在所有可用路径上，部分链路过载而其他链路闲置，造成网络资源的严重浪费。

因此，学术界和工业界亟需重新思考数据中心网络的设计理念，从通用计算优化转向专门针对AI工作负载的定制化架构设计。本章将系统性地分析Meta在这一领域的工程实践，从网络架构的革新、超大规模集群的互联技术，到计算与存储的协同设计，全面展示大规模AI基础设施设计的核心原理与技术实现。

\section{从传统网络到专用AI网络}

\subsection{传统Clos网络的局限性}

传统数据中心网络广泛采用Clos拓扑结构，这种架构通过多级交换机的递归连接实现无阻塞（non-blocking）或低阻塞的网络特性。在通用云计算环境中，Clos架构表现出色，因为其流量模式相对分散，单个流的持续时间较短，且流数量众多，能够充分利用ECMP机制实现负载均衡。

然而，当面对AI训练工作负载时，传统Clos网络暴露出一系列结构性问题。首先，大象流的存在使得基于流的哈希分配策略失效。由于AI训练过程中的梯度同步、参数更新等操作涉及海量数据的持续传输，单个流可能长时间占用某条链路的大部分带宽。当多个大象流通过哈希算法被分配到同一条链路时，该链路过载而其他链路空闲，产生严重的负载不均衡现象——正如我们之前所说的，多条超长货车队挤在同一条车道，而其他车道却空无一车。

其次，低熵流量特征进一步加剧了哈希碰撞问题。在集体通信操作（如AllReduce、AllGather）中，参与通信的GPU数量有限，导致生成的IP流数量相对较少。有限的流数量意味着哈希空间被不充分探索，流与链路之间的映射容易出现冲突。这就像是在一个只有几条出口的高速收费站，大量车辆被迫挤向有限的通道，即使其他通道完全空闲。即使网络整体具备充足的带宽容量，局部拥塞仍会导致训练性能下降。

Meta的工程团队曾尝试多种方法来缓解这些问题。例如，他们设计了基于边界网关协议（BGP）的策略，使从加速器经由叶交换机接收的流量根据目的地固定到特定上行链路。这种方法在稳定状态下缓解了低熵问题，但在故障场景下失效，因为此时回退机制仍依赖ECMP路由。他们还尝试了负载感知的ECMP方案，虽然能够处理胖流和低熵问题，但调参困难且会产生乱序包，这对RDMA通信极为不利——RDMA要求数据包按序到达，乱序会导致严重的性能退化。此外，基于流量工程的解决方案预先计算特定模型使用的流模式，在作业开始前配置叶交换机，但这种方法随着网络规模增长变得过于复杂，且其中央化特性使其对故障的响应迟缓。

\subsection{向专用AI网络的演进}

上述挑战表明，修补传统网络架构无法从根本上解决问题，必须采用全新的设计理念。这一认识推动了专用AI网络架构的诞生，其核心思想是将网络从通用数据包转发系统转变为针对AI工作负载优化的数据调度系统。

专用AI网络的关键特征包括：支持细粒度的负载均衡机制以应对大象流，采用无哈希的转发策略以避免低熵问题，实现无损（lossless）传输以支持RDMA操作，以及提供可扩展的架构以支持数万GPU的互联。这些需求催生了多种创新技术，其中最具代表性的是\textbf{解耦式调度网络（Disaggregated Scheduled Fabric, DSF）}和\textbf{后端聚合（Backend Aggregation, BAG）}技术。

这里需要明确两者的关系与定位：\textbf{DSF和BAG是专用AI网络架构的两个互补组成部分}。DSF负责解决数据中心\textbf{内部}的互联挑战，通过创新的交换架构实现机架间的高效通信；而BAG则解决数据中心\textbf{之间}乃至跨区域的大规模互联问题，作为连接多个DSF区域的"超级脊层"。可以把整个网络想象成一座城市的高速公路系统：DSF是城市内部的道路网络，负责将车辆（数据包）从起点高效送达目的地；BAG则是连接多个城市的高速铁路，负责跨区域的快速运输。两者协同工作，共同构成了支撑超大规模AI训练的完整网络基础设施。下文将首先介绍BAG技术，因为它代表了从宏观视角解决跨区域互联问题的思路。

\section{后端聚合：连接万卡集群}

\subsection{BAG的设计理念}

随着AI集群规模从单数据中心扩展到多数据中心甚至跨区域部署，如何在保持性能的同时实现超大规模互联成为核心挑战。Meta提出的后端聚合（Backend Aggregation, BAG）技术正是为解决这一问题而设计的中心化以太网超级脊层网络。

\textbf{BAG拓扑的核心特征}是采用集中式超级脊层连接多个L2网络结构。具体来说，BAG层位于区域网络与Meta骨干网之间，充当多个数据中心之间的"交通枢纽"。想象一下：多个城市的机场（L2结构）通过一个大型航空枢纽（BAG超级脊层）连接起来，乘客（数据包）可以先到达枢纽，再转机前往最终目的地。这种设计将分布式计算资源通过高容量网络抽象为统一的计算池，从而支持吉瓦级别的AI训练基础设施。

BAG的主要功能是跨多个数据中心和区域互连多个脊层网络。以Meta的Prometheus项目为例，BAG层作为区域网络与Meta骨干网之间的聚合点，使得构建超大AI集群成为可能。

BAG设计面临的关键约束包括距离限制、缓冲区容量和延迟要求。由于BAG-to-BAG连接可能跨越较长的物理距离，这要求使用具备深缓冲区（deep buffer）能力的交换机，以支持如基于优先级的流控制（PFC）等无损拥塞控制协议。深缓冲区提供了充足的头部空间（headroom buffer），确保在网络拥塞时不会丢失数据包，这对RDMA通信至关重要。

\subsection{架构实现}

BAG架构采用模块化硬件设计，使用配备Jericho3（J3）ASIC线卡的机箱式交换机，每块线卡提供多达432个800G端口。这种高密度的端口配置使得BAG能够以PB级带宽（例如每对区域之间16-48 Pbps）支持区域间的海量数据传输需求。中心化的BAG集线器使用更大的机箱以容纳众多分支（spoke）和长距离链路，通过优化的缓冲区利用策略适应不同长度的电缆连接。

在拓扑设计上，BAG层在区域间采用平面（planar）或分散（spread）连接拓扑，选择依据主要是站点规模和光纤可用性。平面拓扑按照平面一对一连接BAG交换机，管理相对简单但将潜在的故障域集中。分散连接拓扑则将链路分布在多个BAG交换机和平面上，增强了路径多样性和系统韧性。

从L2结构到BAG的典型\textbf{超配比（oversubscription ratio）}约为4.5:1。超配比表示下行带宽与上行带宽的比值——例如4.5:1意味着下行总带宽是上行总带宽的4.5倍。这是一种成本与性能之间的折中设计：并非所有GPU都需要同时以满带宽与其他区域的GPU通信，因此可以适当减少上行带宽以降低成本。而BAG-to-BAG的超配比则根据区域需求和链路容量灵活调整。

BAG与下层L2结构的连接方式取决于所采用的网络技术。对于DSF结构，BAG通过每个数据中心的特殊后端边缘机柜（edge PoD）连接到DSF L2区域；对于非调度结构（NSF），每个BAG平面连接到来自所有脊平面的脊训练交换机（STSW），形成有效的4.98:1超配比。

\subsection{路由策略与安全性}

BAG内部路由采用外部边界网关协议（eBGP），并引入链路带宽属性以支持非等价多路径（UCMP）负载均衡。UCMP允许流量根据链路容量按比例分配，而非传统的等量分配，从而更高效地利用异构网络资源。这种机制还提供了强大的故障处理能力，当某条链路失效时，流量能够自动重新分配到其他可用路径。

在安全方面，BAG-to-BAG连接采用MACsec加密，确保跨区域数据传输的机密性和完整性，符合企业网络的安全合规要求。此外，网络设计详细规划了端口条带化（port striping）、IP地址分配方案，并进行了全面的故障域分析，以确保高可用性并最小化故障影响范围。

\subsection{故障恢复机制}

大规模网络中的故障不可避免，因此BAG设计了多层次的故障缓解策略。针对黑洞（blackholing）风险，BAG采用多种技术手段，包括主动排空（draining）受影响的BAG平面，以及条件性路由聚合。这些机制确保在硬件故障或链路中断时，网络能够快速收敛，将服务中断时间降至最低。

故障分析覆盖BAG层、数据大厅级别和电力分配级别，确保从芯片级故障到数据中心级灾难的全方位容错能力。这种分层故障管理策略对于支持吉瓦级AI集群的持续运行至关重要。

\section{解耦式调度网络：DSF架构}

\subsection{两域架构设计}

解耦式调度网络（Disaggregated Scheduled Fabric, DSF）代表了数据中心网络架构的根本性创新。其核心思想是打破传统单体机箱式交换机的物理限制，将线卡（line cards）和交换网板（fabric cards）解耦为独立互联的硬件设备，构建分布式交换系统。

\textbf{DSF采用独特的双域架构}，将网络划分为\textbf{以太网域}和\textbf{交换结构域}。这两个域就好比邮政系统中的不同环节：以太网域是传统的街道地址系统，负责基于标准以太网协议的数据包路由；交换结构域则像是现代化的包裹分拣中心，负责高效地将"包裹"（数据包单元）快速分发到正确的出口。

具体来说，在以太网域中，服务器和传统网络协议正常运行；而在交换结构域中，数据包被分割为单元（cells），通过包喷洒（packet spraying）技术均匀分布在所有可用路径上，在出接口处重新组装后交付给以太网域。这种设计的精妙之处在于，对外部网络基础设施而言，分布式的接口节点和交换节点表现为单一统一的交换机，其外部端口总数等于所有接口节点外部端口的总和，有效创建了超越物理限制的虚拟机箱交换机。

DSF架构由两类核心组件构成：\textbf{接口节点（Interface Nodes, INs）}，也称为机架解耦交换机（RDSWs）；以及\textbf{交换节点（Fabric Nodes, FNs）}，即交换网解耦交换机（FDSWs）。接口节点负责外部连接和路由功能，与更广泛的数据中心基础设施交互；交换节点则作为内部交换单元，专注于高速流量在结构内部的分发，无需三层路由能力。

\subsection{控制平面与操作系统}

DSF的控制平面基于Meta开源的网络操作系统FBOSS构建，支持解耦式结构的多ASIC控制需求。FBOSS通过与FBOSS状态数据库（FSBD）通信，实现跨节点的实时状态同步。这种设计确保在分布式环境下，网络状态的一致性和及时更新，为动态流量调度提供基础。

FBOSS的开放性和可编程性使DSF能够快速适应不断演进的AI工作负载需求，同时也为网络故障诊断和性能优化提供了丰富的工具和接口。

\subsection{虚拟输出队列机制}

\textbf{虚拟输出队列（Virtual Output Queuing, VOQ）}是DSF实现无损传输的核心机制。可以把VOQ想象成一个现代化的机场值机系统：传统机场只有一个统一的等待区（输入队列），如果最前面的乘客遇到了问题，后面所有乘客都被堵住——这就是\textbf{队头阻塞（head-of-line blocking）}问题。而VOQ就像是机场为每个目的地航班设立了独立的候机室，每个目的地的乘客在各自的候机室等待，互不影响。

具体来说，\textbf{VOQ机制为每个目标端口独立维护虚拟队列}。在DSF中，进入的数据包根据其目标端口被导向相应的虚拟输出队列，每个队列独立调度传输。这种细粒度的流量管理能够精确适应AI工作负载的通信模式，避免不同类型流量之间的相互干扰。当某个输出端口拥塞时，只有指向该端口的流量受到影响，其他流量继续正常传输，从而最大化网络吞吐量。

\subsection{包喷洒技术}

包喷洒是DSF最具创新性的特性之一，也是其区别于传统以太网结构的关键所在。与依赖哈希的传统方法不同，DSF利用硬件能力在结构域内将数据包分割为单元并均匀喷洒在所有可用路径上，同时在接口节点处确保有序交付给端主机。

这种技术的实现依赖于基于信用的分配机制（credit-based allocation scheme）。入口接口节点动态地向出口接口节点请求信用令牌，系统根据当前路径可用性、拥塞水平和带宽利用率做出实时决策。这种机制使DSF能够在所有可用网络路径上实现接近最优的负载均衡，有效利用结构的全部带宽容量。

包喷洒技术的优势在处理混合流量模式和动态网络条件时尤为明显。无论流量特征如何变化，DSF都能自动适应，无需人工重新配置或流量工程干预。这与传统ECMP形成鲜明对比——后者在流量模式变化时往往需要进行复杂的重新调优。

\subsection{DSF拓扑扩展}

DSF架构通过多级扩展支持不同规模的AI集群。在基础层面，AI区域（AI zone）包含多个扩展单元（scaling units），每个单元是一组连接到该区域RDSW的GPU机架。同一AI区域内的所有RDSW通过公共的FDSW层互联。RDSW采用深缓冲区Jericho3-AI芯片，FDSW则使用Ramon3芯片，两者均运行FBOSS操作系统。RDSW与FDSW之间使用2x400G FR4光模块连接。

为支持单一AI区域内的更高GPU规模，DSF创建两个相同的网络平面，构成DSF L1区域，这是更大规模GenAI集群的构建模块。通过第二级脊DSF交换机（SDSW）互连四个DSF L1区域，形成DSF L2区域。SDSW使用与FDSW相同的硬件，通过非阻塞拓扑实现高达18,000块800G GPU的互联规模。

在更大规模部署中，五个DSF L2区域通过L3超级脊层互联。每个数据中心的边缘机柜（Edge PoD）包含40个FDSW和128个边缘DSF交换机（EDSW）。EDSW在硬件上与RDSW相同，但功能上专用于提供到L3超级脊的连接。每个EDSW使用4x800G链路连接到四个超级脊设备，每个边缘机柜总计提供2,000个800G端口。

在L3互联层面，路由信息的交换采用内部BGP（iBGP）会话连接EDSW与建筑物内所有RDSW，启用BGP add-path功能使RDSW能够通过全部2,000个下一跳学习聚合路由。EDSW与L3超级脊之间使用eBGP，仅交换聚合路由信息。由于训练模型的分片特性，L3超级脊层的流量预期较少，因此4.5:1的超配比足以满足需求。

\subsection{输入平衡模式}

DSF引入了一项创新特性——输入平衡模式（input-balanced mode）。该功能智能地平衡各层之间的可达性信息，确保在发生故障时避免在结构层和脊层出现拥塞。这种前瞻性的设计体现了DSF对大规模网络中故障常态化的深刻认识，通过协议层面的优化提升系统整体韧性。

\section{存储与计算协同设计}

\subsection{Tectonic分布式存储系统}

大规模AI训练不仅对网络架构提出挑战，也对存储系统提出了严苛要求。随着训练任务日益多模态化，消耗大量图像、视频和文本数据，数据存储需求呈指数级增长。同时，将这些海量存储整合到高性能、高能效的架构中成为必须解决的难题。

Meta开发了Tectonic分布式存储解决方案来满足AI集群的数据和检查点（checkpointing）需求。Tectonic采用基于Linux用户态文件系统（FUSE）的API，针对闪存介质进行优化，使数千块GPU能够以同步方式保存和加载检查点。检查点同步是大规模分布式训练中的关键挑战——任何延迟都会导致昂贵的计算资源闲置。

在硬件层面，存储部署基于YV3 Sierra Point服务器平台，配备市场上最高容量的E1.S SSD。除了更高容量的SSD外，每机架的服务器数量经过定制，以实现每台服务器的吞吐量容量、机架数量减少和相关能效之间的最佳平衡。利用OCP服务器如同乐高积木般的构建块特性，存储层能够灵活扩展以满足当前和未来更大规模AI集群的需求，同时具备对日常基础设施维护操作的容错能力。

\subsection{Hammerspace并行文件系统}

除Tectonic外，Meta还与Hammerspace合作开发并部署了并行网络文件系统（NFS），以满足AI集群的开发者体验需求。Hammerspace的一项关键优势是使工程师能够对使用数千块GPU的作业执行交互式调试——代码变更能够立即在环境中的所有节点上可见。

\textbf{Tectonic和Hammerspace构成了分层存储架构的两个互补层次}。可以把这种架构类比为计算机的内存层次结构：Tectonic相当于高速缓存（Cache），专注于为AI训练提供高吞吐量的数据存取和检查点服务；而Hammerspace则相当于主内存，优化开发者的交互式工作体验。两者的区别在于：\textbf{Tectonic专注于训练数据的高吞吐存取}，针对大规模并行读写进行了深度优化；\textbf{而Hammerspace优化开发者的交互式调试体验}，提供类POSIX文件语义和即时全局可见性。

Tectonic分布式存储与Hammerspace的结合实现了快速迭代速度，同时不牺牲规模。这种分层存储架构设计反映了Meta对AI研发工作流程的深刻理解：高性能的分布式存储满足训练吞吐需求，而并行文件系统则优化了开发者的交互体验。

\subsection{Grand Teton计算平台}

计算层面，Meta采用自主设计的Grand Teton开放GPU硬件平台。Grand Teton是在多年AI系统集成经验基础上发展的产物，将电源、控制、计算和结构接口集成到单一机箱中，实现更优的整体性能、信号完整性和热性能。它提供快速的可扩展性和简化的设计，能够快速部署到数据中心集群并易于维护和扩展。

Grand Teton与Meta的其他内部创新（如Open Rack电源和机架架构）相结合，使Meta能够以专为当前和未来应用定制的方式构建新集群。这一设计哲学体现了软硬件协同优化的思想——计算平台的设计充分考虑了网络结构和存储系统的特性，形成端到端的优化方案。

Meta自2015年的Big Sur平台起就开始开放设计其GPU硬件平台，这一传统延续到Grand Teton，并贡献给开放计算项目（OCP），推动行业开放创新。

\section{性能优化与调试}

\subsection{大规模集群的性能挑战}

构建大规模AI集群的难点不仅在于硬件部署，更在于充分释放硬件潜力。Meta的实践经验表明，即使拥有顶尖的网络设备和计算资源，未经优化的大规模集群性能可能远低于预期。具体而言，小型集群的通信带宽和利用率能够达到90\%以上，但未优化的大型集群利用率可能从10\%到90\%大幅波动，性能极不稳定。

这种性能退化的根本原因在于规模扩大带来的复杂性。随着GPU数量增加，网络拓扑变得更加复杂，流量模式更加难以预测，故障概率相应提高。传统的针对小规模集群的优化策略在大规模环境中失效，需要全新的系统级优化方法。

\subsection{拓扑感知调度}

为应对上述挑战，Meta对其内部作业调度器进行了重大改进，引入网络拓扑感知能力。\textbf{拓扑感知调度的核心思想可以概括为"把经常聊天的GPU放在同一间宿舍"}——确保相互通信频繁的GPU尽可能部署在物理距离较近的位置，通常是同一机架或同一叶交换机下。

这种策略带来了两方面收益：首先减少了通信延迟，因为数据无需经过多级交换机；其次最小化了向上层网络传输的流量，减轻了脊层和超级脊层的负载压力。可以设想一个简单例子：假设有4个GPU，A和B之间、C和D之间通信频繁，但A与C、B与D之间几乎没有通信。拓扑感知调度会将A和B部署在同一机架，C和D部署在另一机架，这样机架间只需传输极少的流量。

拓扑感知调度需要在资源利用率和通信性能之间取得平衡。过于严格的局部性要求可能导致资源碎片化，降低整体利用率；而过于宽松的策略则无法获得网络局部性带来的性能优势。Meta的解决方案体现了对这一权衡的深刻把握。

\subsection{网络路由优化}

除调度优化外，Meta还优化了网络路由策略，并与NVIDIA集合通信库（NCCL）协同改进以实现最佳网络利用率。NCCL是GPU间高效通信的基础软件，其配置直接影响网络流量的分布特征。通过与NCCL的协同优化，Meta确保上层应用生成的流量模式与下层网络架构的最佳工作点相匹配。

这种软硬件协同优化是大规模AI集群设计的核心原则之一。单独优化网络硬件或软件都无法达到最佳效果，只有将两者作为一个整体系统设计，才能充分释放规模带来的算力红利。

\subsection{故障诊断工具}

可调试性被Meta识别为大规模训练中的主要挑战之一。当训练作业因某块GPU停滞而中断时，在数千块GPU中定位问题根源极其困难。为解决这一问题，Meta正在开发诸如desync debug和分布式集合通信飞行记录器（distributed collective flight recorder）等工具。

这些工具暴露分布式训练的细节，帮助工程师更快、更容易地识别问题。desync debug能够检测并定位同步点处的性能异常，而飞行记录器则捕获集合通信的历史状态，为事后分析提供数据支持。这些诊断能力的建设对于运维大规模AI集群至关重要，因为随着规模扩大，故障从罕见事件转变为常态，快速诊断和恢复能力直接影响集群的有效利用率。

\subsection{框架层优化}

在PyTorch框架层面，Meta识别并解决了多个性能瓶颈。其中一个关键改进针对进程组初始化过程——在大规模集群中，这一过程的启动时间可能长达数小时。通过系统性优化，Meta将启动时间缩短至分钟级别，显著提升了开发者体验和集群周转效率。

此外，充分利用NVIDIA H100 GPU支持的新数据类型（如8位浮点数FP8）进行训练，需要投资额外的并行化技术。Meta与训练框架和模型编写团队紧密合作，使软件栈能够适应不断演进的基础设施，充分利用硬件新特性带来的性能增益。

\subsection{存储检查点优化}

在存储层面，Meta针对数千个训练rank的检查点操作进行了深度优化，使其能够在数百毫秒内完成。这一优化对于提升大规模训练的有效时间占比至关重要——如果检查点操作耗时过长，昂贵的GPU计算资源将被迫等待，造成巨大浪费。

\section{小结}

本章系统性地分析了Meta在大规模AI基础设施设计方面的工程实践，从网络架构的革新到计算存储的协同优化，展现了构建吉瓦级AI集群的完整技术图景。以下是本章的核心要点：

\begin{enumerate}
    \item \textbf{网络架构革新}：传统Clos网络因大象流和低熵流量问题不再适用于AI训练，这推动了专用AI网络的诞生。大象流就像高速公路上的超长货车队，低熵问题就像少数车道承载大量车辆，都会导致严重的负载不均衡。
    
    \item \textbf{BAG与DSF的协同}：\textbf{后端聚合（BAG）}作为"超级脊层"负责跨区域互联，采用集中式超级脊层连接多个L2结构；\textbf{解耦式调度网络（DSF）}则负责数据中心内部互联，通过双域架构（以太网域+交换结构域）实现高效通信。两者共同构成了完整的AI网络基础设施。
    
    \item \textbf{DSF的关键创新}：DSF通过\textbf{虚拟输出队列（VOQ）}机制——为每个目标端口独立维护虚拟队列——解决了队头阻塞问题；通过\textbf{包喷洒技术}实现无损负载均衡，突破了物理交换机的限制。
    
    \item \textbf{分层存储架构}：\textbf{Tectonic}专注于训练数据的高吞吐存取，\textbf{Hammerspace}优化开发者的交互式调试体验，两者形成互补的分层存储架构。
    
    \item \textbf{系统性优化}：从拓扑感知调度（把经常通信的GPU放在同一间宿舍）到框架层优化，从网络路由策略到故障诊断工具，系统性优化是充分释放规模效益的必要条件。
\end{enumerate}

展望未来，随着AI模型规模和数据量的持续增长，基础设施设计将面临更大挑战。Meta的实践表明，开放创新、软硬件协同、以及针对工作负载特性的定制化设计，是应对这些挑战的有效路径。这些经验不仅适用于Meta的特定场景，也为整个行业构建下一代AI基础设施提供了宝贵参考。
